% ============================================================================
% APPENDIX
% ============================================================================
\chapter{Source Code Listings}

\section{Research Agent --- Core Tournament Function}

\begin{lstlisting}[language=Python, caption={run\_research\_agent() --- Tournament-Based Execution}, label={lst:app_research}]
def run_research_agent(topic, domain, role, years=None):
    """Spawns multiple agents with varying temperatures,
    evaluates all outputs, selects highest scoring."""
    base_prompt = build_prompt(topic, domain, role, years)
    num_parallel = 3

    def _run_single_experiment(experiment_id):
        temp = 0.2 + (experiment_id * 0.2)
        exp_prompt = base_prompt + (
            "\n\nFINAL REQUIREMENT: Return one JSON object."
        )
        if experiment_id == 1:
            exp_prompt += "\nFocus on technical architectures."
        elif experiment_id == 2:
            exp_prompt += "\nFocus on unsolved edge-cases."
        try:
            raw = call_llm(exp_prompt, role, temperature=temp)
            parsed = extract_json_only(raw)
            score, feedback = evaluate_research_output(
                parsed, topic, role)
            return {"score": score, "parsed": parsed,
                    "feedback": feedback, "error": None}
        except Exception as e:
            return {"score": -1, "parsed": None,
                    "error": str(e)}

    with ThreadPoolExecutor(max_workers=num_parallel) as ex:
        futures = [ex.submit(_run_single_experiment, i)
                   for i in range(num_parallel)]
        results = [f.result()
                   for f in as_completed(futures)]
    results.sort(key=lambda x: x["score"], reverse=True)
    best = results[0]
    if best["score"] == -1:
        return {"status": "failed", "error": best["error"]}
    final = _ensure_defaults(best["parsed"],
                             topic, domain, role)
    final["_autoresearch_score"] = best["score"]
    return final
\end{lstlisting}

\section{Streamlit Configuration}

\begin{lstlisting}[caption={.streamlit/config.toml --- Server Configuration}, label={lst:config}]
[server]
scriptRunTimeout = 0
enableWebsocketCompression = false
maxUploadSize = 200
maxMessageSize = 200
enableXsrfProtection = false
headless = true
enableStaticServing = true

[browser]
gatherUsageStats = false
serverAddress = "localhost"

[runner]
fastReroute = true
magicEnabled = false
\end{lstlisting}

\section{Environment Configuration Template}

\begin{lstlisting}[caption={.env.example --- API Key Configuration}, label={lst:env}]
GROQ_API_KEY=gsk_your_api_key_here
GROQ_MODEL=llama3-70b-8192
OLLAMA_URL=http://localhost:11434
OLLAMA_MODEL=mistral
\end{lstlisting}

\section{Orchestrator Utilities}

\begin{lstlisting}[language=Python, caption={orchestrator/utils.py --- Shared Utilities}, label={lst:utils}]
import json, uuid
from datetime import datetime, timezone

def now_iso_z():
    return datetime.now(timezone.utc).replace(
        microsecond=0).isoformat().replace(
        "+00:00", "Z")

def make_agent_id(prefix, role):
    return f"{prefix}_{role}_{uuid.uuid4().hex[:8]}"

def load_prompt(path):
    with open(path, "r", encoding="utf-8") as fh:
        return fh.read()
\end{lstlisting}

\section{Prompt Template Renderer}

\begin{lstlisting}[language=Python, caption={agents/prompt\_helpers.py --- Template Engine}, label={lst:prompt}]
def simple_render(template_text, mapping):
    out = template_text
    for k, v in mapping.items():
        token = "{{" + k + "}}"
        out = out.replace(token, str(v))
    return out
\end{lstlisting}
